% Appendix fragment; requires amsmath and amssymb only.
\section{Why atomic case introduction cannot simply be removed}
\label{app:forced-progress-obstruction}

\paragraph{Attribution.} The counterinterpretation and the lost theorem $D_{x,y}(y=1)$ already appear in the current flat-computation draft. The following is a self-contained rederivation of that inherited obstruction, not a new contribution.~\cite{flat}

Another tempting reduction is to make the introduction mode automatic:
use progress when the matrix is a literal equation, and case introduction
when it is compound.  This would remove one certificate choice.  It does
not preserve the native theory with its existing elimination rule.

Recall that $\mathsf Q_yB(y)$ is an internal clause with meaning
$B(0)\to\forall yB(y)$, and that paired elimination of
$\mathsf D_{x,y}B(x,y)$ supplies such a clause at a fresh witness.
The unrestricted introduction inventory allows
\[
 \frac{\Gamma,B(t,0)\vdash B(t,Sz)}
      {\Gamma\vdash\mathsf D_{x,y}B(x,y)}
\]
for every matrix, and
\[
 \frac{\Gamma,E(t,z)\vdash E(t,Sz)}
      {\Gamma\vdash\mathsf D_{x,y}E(x,y)}
\]
for literal equation matrices.  Both rules have the usual freshness
conditions.  The proposed reduction deletes the first rule exactly when
$B$ is a literal equation.

\paragraph{A sound alternative semantics for the reduced rules.}
Take the standard natural numbers for all numerical terms and actual
equality for equations.  Interpret $\mathsf Q$ as above.  Interpret
$\mathsf D$ recursively by
\[
 (\mathsf D_{x,y}B)^\dagger=
 \begin{cases}
 \exists x\,\forall y\,
   \bigl(B^\dagger(x,y)\to B^\dagger(x,Sy)\bigr),
       & B\text{ is a literal equation},\\[2pt]
 \exists x\,
   \bigl(B^\dagger(x,0)\to\forall yB^\dagger(x,y)\bigr),
       & B\text{ is compound}.
 \end{cases}
\]
This is a definition on each fixed finite formula, not a truth predicate
inside the calculus.

Forced progress introduction is sound by its direct meaning.
Compound case introduction is sound by the zero-or-successor property
of the standard natural numbers.  For paired elimination with an atomic
matrix, ordinary numerical induction turns the supplied progressiveness
into its conditional universal closure.  Paired elimination with a
compound matrix and closure specialization are sound directly.  Equality
transport is sound because numerical substitution leaves the distinction
between a literal equation and a compound matrix unchanged.  Numerical
computation and every context-availability operation are sound as well.
Therefore every theorem of the reduced rule system is true under
$\dagger$.

\paragraph{A lost theorem.}
Let $1=S0$ and let the witness coordinate be unused.  The original
calculus proves
\[
                   \mathsf D_{x,y}(y=1).
\]
Indeed, case introduction assumes its zero instance $0=1$; numerical
explosion yields $Sz=1$, so the rule applies.  Numerical explosion here
only uses the existing numerical selector and equality transport: an
assumption $0=1$ identifies its two endpoints and hence yields any
literal equation.  No induction on a compound formula is required.

Under $\dagger$, however, this displayed formula says
\[
                   \exists x\,\forall y\,(y=1\to Sy=1),
\]
which is false at $y=1$.  Consequently it is not derivable in the reduced
calculus.

The padding equivalence between $A$ and a suitably gated paired closure
does not repair this counterexample.  Padding can move a proof obligation
into a compound matrix, but the existing elimination interface only
returns conditional closure; it does not recover the progressiveness
required to introduce an arbitrary raw atomic paired formula.
Thus removing atomic case introduction would require a substantive
change of rules and a new interpretation-composite proof.  It is not an
admissible simplification of the present kernel.
